
\chapter{Centralización de vulnerabilidades mediante DefectDojo}

\section{Descripción de DefectDojo}

DefectDojo es una plataforma de código abierto para la gestión de
vulnerabilidades. Recibe los reportes de distintas herramientas de seguridad,
los convierte a un formato común y los organiza en una jerarquía de
\textit{producto}, \textit{engagement} y \textit{test}. Sobre esa base permite
clasificar los hallazgos por severidad, deduplicarlos entre ejecuciones,
asignarles un estado (activo, verificado, mitigado, falso positivo) y dar
seguimiento a su corrección.

En este proyecto cumple una función concreta: los resultados de Semgrep,
Trivy, Gitleaks, TruffleHog y OWASP ZAP tienen formatos distintos y viven en
archivos y \textit{artifacts} separados. DefectDojo los reúne en un solo
lugar, de modo que el equipo consulta un único panel en lugar de abrir cada
reporte. Se utilizó una instancia local desplegada con Docker Compose, y el
producto creado se llama \texttt{MoonPhases}.

\section{Integración con las herramientas de seguridad}

\subsection{Procedimiento automático}

La carga automática la realiza el script
\texttt{scripts/upload\_to\_defectdojo.sh}, invocado por el job
\texttt{upload-defectdojo} del pipeline de GitHub Actions. El job se ejecuta
al final, después de los escaneos, descarga los reportes publicados como
\textit{artifacts} y envía cada uno a la API REST de DefectDojo.

\begin{table}[H]
\centering
\caption{Reportes importados a DefectDojo y tipo de escaneo.}
\label{tab:dojo-reportes}

\begin{tabularx}{\textwidth}{l l X}
\toprule
\textbf{Herramienta} &
\textbf{Tipo de escaneo (parser)} &
\textbf{Reporte} \\
\midrule
Semgrep & \texttt{Semgrep JSON Report} &
\texttt{reports/sast/semgrep-report.json} \\
Trivy & \texttt{Trivy Scan} &
\texttt{reports/sca/trivy-report.json} \\
Gitleaks & \texttt{Gitleaks Scan} &
\texttt{reports/secrets/gitleaks-report.json} \\
TruffleHog & \texttt{Trufflehog Scan} &
\texttt{reports/secrets/trufflehog-report.json} \\
OWASP ZAP & \texttt{ZAP Scan} &
\texttt{reports/dast/baseline.xml} y \texttt{baseline-api.xml} \\
\bottomrule
\end{tabularx}
\end{table}

El script usa el servicio \texttt{reimport-scan} de la API en lugar de una
importación simple, con un \textit{engagement} fijo por rama
(\texttt{CI - <rama>}) y un \textit{test} fijo por reporte. Así DefectDojo
reconoce los hallazgos que se repiten entre ejecuciones y no los duplica en
cada \textit{pipeline}; además, con \texttt{close\_old\_findings} marca como
cerrados los hallazgos que dejaron de aparecer. Cada importación registra
también la rama, el identificador de ejecución y el \textit{commit}.

Reglas de funcionamiento del script:

\begin{itemize}
    \item Si un reporte está vacío (Gitleaks y TruffleHog sin hallazgos), se
    omite sin error.
    \item Si un reporte no existe (por ejemplo, los de ZAP en el
    \textit{runner}), se omite y se informa.
    \item Si una importación falla, el script continúa con los demás reportes
    y termina con error al final.
    \item La URL, el token y el identificador del producto se leen de
    variables y secretos del repositorio de GitHub; el token nunca se escribe
    en el código.
\end{itemize}

\begin{lstlisting}[language=yaml, caption={Job que sube los hallazgos a DefectDojo.}, label={lst:dojo-job}]
upload-defectdojo:
  needs: [gitleaks-scan, trufflehog-scan, trivy-scan, semgrep-scan]
  if: ${{ always() && vars.DEFECTDOJO_URL != '' }}
  env:
    DEFECTDOJO_URL: ${{ vars.DEFECTDOJO_URL }}
    DEFECTDOJO_PRODUCT_ID: ${{ vars.DEFECTDOJO_PRODUCT_ID }}
    DEFECTDOJO_TOKEN: ${{ secrets.DEFECTDOJO_TOKEN }}
  steps:
    - uses: actions/checkout@v4
    - uses: actions/download-artifact@v4   # uno por reporte
    - run: bash scripts/upload_to_defectdojo.sh
\end{lstlisting}

\subsection{Procedimiento manual y limitaciones}

El mismo script permite subir un reporte de forma manual indicando el tipo de
escaneo y el archivo (\texttt{./upload\_to\_defectdojo.sh <tipo>
<archivo>}), y DefectDojo admite además la carga desde su interfaz web, como
se practicó en la guía de laboratorio del curso. Este procedimiento es el
adecuado para los reportes de ZAP, que se generan en local y no están en el
\textit{pipeline}. Su carga con el script está prevista, pero no se ha
ejecutado todavía.

Existe además una limitación de infraestructura: la instancia de DefectDojo
corre en un equipo local y el \textit{runner} de GitHub Actions está en la
nube, por lo que no puede alcanzar \texttt{localhost}. Para las pruebas se
expuso DefectDojo mediante un túnel de ngrok, cuya dirección cambia si se
reinicia el túnel. En un entorno estable haría falta un servidor accesible
públicamente o un \textit{runner} propio.

\evidencefigure
{figures/DefectDojo.png}
{Panel de DefectDojo tras la importación de los reportes del pipeline.}
{fig:defectdojo}

\section{Flujo de gestión de vulnerabilidades}

El proceso completo, desde que una herramienta genera un reporte hasta que el
hallazgo se mitiga, sigue estos pasos:

\begin{enumerate}
    \item \textbf{Generación:} cada herramienta de seguridad produce su
    reporte en JSON o XML durante el \textit{pipeline}.
    \item \textbf{Importación:} el job \texttt{upload-defectdojo} envía los
    reportes a DefectDojo. Los hallazgos se deduplican y quedan asociados al
    producto, a la rama y a la ejecución que los originó.
    \item \textbf{Clasificación:} el equipo revisa los hallazgos nuevos,
    confirma su severidad y descarta los falsos positivos. Cada hallazgo se
    relaciona con una amenaza del modelo STRIDE o se registra como nueva.
    \item \textbf{Seguimiento:} se asigna un responsable y se prioriza según
    severidad e impacto (capítulo siguiente).
    \item \textbf{Mitigación:} se corrige el hallazgo en el código o en la
    configuración y se vuelve a ejecutar el \textit{pipeline}.
    \item \textbf{Verificación:} en la siguiente importación, un hallazgo
    corregido deja de aparecer y DefectDojo lo marca como cerrado, lo que
    deja registro de su ciclo de vida.
\end{enumerate}

\begin{infobox}
\textbf{Flujo general:}

Herramientas de seguridad $\rightarrow$
Reportes $\rightarrow$
Importación $\rightarrow$
DefectDojo $\rightarrow$
Clasificación $\rightarrow$
Mitigación.
\end{infobox}

\section{Resultados de la centralización}

En la ejecución del pipeline del 23 de noviembre de 2026, el job
\texttt{upload-defectdojo} terminó correctamente. El registro de la
ejecución muestra el resultado de cada reporte:

\begin{table}[H]
\centering
\caption{Resultado de la importación en la ejecución del pipeline.}
\label{tab:dojo-resultado}

\begin{tabularx}{\textwidth}{l c X}
\toprule
\textbf{Reporte} &
\textbf{Hallazgos} &
\textbf{Resultado de la importación} \\
\midrule
Semgrep & 15 & Importado (HTTP 201). \\
Trivy & 1 & Importado (HTTP 201). \\
Gitleaks & 0 & Omitido: reporte vacío (sin secretos). \\
TruffleHog & 0 & Omitido: reporte vacío (sin secretos). \\
OWASP ZAP & -- & No ejecutado en el \textit{pipeline}. \\
\midrule
\textbf{Total importado} & \textbf{16} & \\
\bottomrule
\end{tabularx}
\end{table}

Los 16 hallazgos importados se distribuyen así según los reportes de origen:
15 advertencias de Semgrep (todas \textit{WARNING}, por acciones de GitHub sin
fijar por \textit{hash}) y 1 vulnerabilidad de severidad alta de Trivy
(\texttt{body-parser}). Los dos reportes de secretos no tienen hallazgos, y
los de ZAP (13 alertas en la aplicación web y 17 en la API) esperan su carga
manual.

La figura~\ref{fig:defectdojo} muestra el panel de DefectDojo después de la
importación. Coincide con los totales de la tabla anterior:

\begin{itemize}
    \item \textbf{16 hallazgos} registrados en los últimos siete días, con
    2 \textit{engagements} activos.
    \item \textbf{Distribución por severidad:} el gráfico de hallazgos
    reportados muestra 15 de severidad media y 1 de severidad alta. Los 15
    medios corresponden a las advertencias de Semgrep y el alto a la
    vulnerabilidad de \texttt{body-parser} detectada por Trivy. No hay
    hallazgos críticos ni bajos.
    \item \textbf{Estado de seguimiento:} 0 hallazgos cerrados y 0 con
    riesgo aceptado. Todos continúan activos, lo que concuerda con que
    ninguno ha sido corregido todavía (capítulo siguiente).
\end{itemize}
